\section{Conclusion}
\label{sec:conclusion}

We revisited anisotropy in large language model training and uncovered a structural asymmetry:
\emph{activation anisotropy is dominated by a coherent mean bias.}
This mean component emerges naturally from hierarchical semantic aggregation
and propagates through residual pathways,
while gradient outputs remain largely mean-neutral.

In high dimensions, such coherent bias amplifies elementwise outliers
and inflates adaptive second-moment statistics,
becoming the primary driver of numerical instability.
We show that this dominant component can be removed
with a simple rank-one mean subtraction,
optionally followed by lightweight spike whitening,
providing effective spectral conditioning at minimal computational cost.

These findings reposition anisotropy as both a \emph{curse} and a \emph{blessing}:
the same structured bias that destabilizes training
also enables efficient mitigation when isolated at its source.